\documentclass[12pt,a4paper]{article}
\usepackage[T1]{fontenc}
\usepackage[utf8]{inputenc}
\usepackage[turkish]{babel}
\usepackage{amsmath,amssymb}
\usepackage{siunitx}
\usepackage{geometry}
\usepackage{physics}
\usepackage{bm}
\geometry{margin=2.3cm}
\sisetup{locale=DE,per-mode=symbol}

\title{Bulut Fiziği ve Yağış Dersi -- Örnek Soru Çözümleri}
\author{Detaylı Çözüm Notu}
\date{}

\begin{document}
\maketitle

Bu dokümanda verilen örnek sorular, ara sınav hazırlığı için detaylı adımlarla çözümlenmiştir.

%%%%%%%%%%%%%%%%%%%%%%%%%%%%%%%%%%%%%%%%%%%%%%%%%%%%%%%%%%%%
\section*{Soru 1 -- Yoğunlaşma Isısı ile Isınma}

Verilenler:
\begin{itemize}
  \item Dikey hız: \(w = \SI{20}{\meter\per\second}\)
  \item Yoğuşan nem: \(\Delta q_v = 0.005\)
  \item Gizli ısı: \(L = 2.5 \times 10^6\)
  \item \(c_p = 1004\)
\end{itemize}

Birim kütlede açığa çıkan ısı:
\[
\Delta Q = L \Delta q_v = 1.25\times 10^4 \, \text{J/kg}
\]

Sıcaklık artışı:
\[
\Delta T = \frac{\Delta Q}{c_p} \approx 12.4 \, \text{K}
\]

Dikey tabakayı geçiş süresi:
\[
\Delta t = \frac{1000}{20} = 50 \, \text{s}
\]

Isınma oranı:
\[
\frac{\Delta T}{\Delta t} \approx 0.25 \,\text{K/s}
\]

%%%%%%%%%%%%%%%%%%%%%%%%%%%%%%%%%%%%%%%%%%%%%%%%%%%%%%%%%%%%
\section*{Soru 2 -- Buharlaşma ve Negatif Yüzdürme}

\[
\Delta T = -\frac{L \Delta q_v}{c_p}
\]

Sayısal:
\[
\Delta T \approx -20 \, \text{K}
\]

Yüzdürme:
\[
B = g \frac{\Delta T}{T} \approx -0.65\, \text{m/s}^2
\]

%%%%%%%%%%%%%%%%%%%%%%%%%%%%%%%%%%%%%%%%%%%%%%%%%%%%%%%%%%%%
\section*{Soru 3 -- Neden \(\Gamma_m < \Gamma_d\)?}

Kuru adyabatik:
\[
\Gamma_d = \frac{g}{c_p} \approx 9.8 \, \text{K/km}
\]

Nemli adyabatik:
\[
\Gamma_m < \Gamma_d
\]

Çünkü yoğuşma gizli ısı açığa çıkarır → soğuma hızı azalır.

%%%%%%%%%%%%%%%%%%%%%%%%%%%%%%%%%%%%%%%%%%%%%%%%%%%%%%%%%%%%
\section*{Soru 4 -- Sanal Sıcaklık Düzeltmeli CAPE}

CAPE:
\[
\text{CAPE} = \int g \frac{T_{v,p} - T_{v,e}}{T_{v,e}} dz
\]

Yaklaşık yüzdürme:
\[
B \approx g \frac{8}{265} \approx 0.30 \, \text{m/s}^2
\]

CAPE:
\[
\text{CAPE} \approx 0.30 \times 4200 \approx 1300 \, \text{J/kg}
\]

Sanal sıcaklık düzeltmesi:
\[
\text{CAPE}_v \approx 1100 \, \text{J/kg}
\]

%%%%%%%%%%%%%%%%%%%%%%%%%%%%%%%%%%%%%%%%%%%%%%%%%%%%%%%%%%%%
\section*{Soru 5 -- CIN Hesabı}

\[
\text{CIN} = -\int g \frac{T_{v,p}-T_{v,e}}{T_{v,e}} dz
\]

Örnek:
\[
\text{CIN} \approx 45\, \text{J/kg}
\]

%%%%%%%%%%%%%%%%%%%%%%%%%%%%%%%%%%%%%%%%%%%%%%%%%%%%%%%%%%%%
\section*{Soru 6 -- Süpersatürasyon Zamanskalası}

Yükselme zamanı:
\[
\tau_{asc} = 40\, \text{s}
\]

Süpersatürasyon:
\[
S \approx \frac{\tau_{asc}}{\tau_{cond}} = 0.2 = 20\%
\]

%%%%%%%%%%%%%%%%%%%%%%%%%%%%%%%%%%%%%%%%%%%%%%%%%%%%%%%%%%%%
\section*{Soru 7 -- Boussinesq Yaklaşımı}

Ana varsayımlar:
\begin{itemize}
  \item \(|\rho'| \ll \rho_0\)
  \item \(\nabla \cdot \mathbf{v} = 0\)
  \item Akustik dalgalar filtrelenir
\end{itemize}

Dikey momentum:
\[
\frac{Dw}{Dt} = -\frac{1}{\rho_0} \frac{\partial p'}{\partial z} + B
\]

%%%%%%%%%%%%%%%%%%%%%%%%%%%%%%%%%%%%%%%%%%%%%%%%%%%%%%%%%%%%
\section*{Soru 8 -- Nemli Adyabatik Lapse Rate Türevi}

\[
\Gamma_m = \Gamma_d 
\frac{1 + \frac{L q_s}{R_d T}}
     {1 + \frac{L^2 q_s}{c_p R_v T^2}}
\]

%%%%%%%%%%%%%%%%%%%%%%%%%%%%%%%%%%%%%%%%%%%%%%%%%%%%%%%%%%%%
\section*{Soru 9 -- CAPE'den Max Updraft}

\[
w_{\max} = \sqrt{2\,\text{CAPE}} = 77.5\, \text{m/s}
\]

%%%%%%%%%%%%%%%%%%%%%%%%%%%%%%%%%%%%%%%%%%%%%%%%%%%%%%%%%%%%
\section*{Soru 10 -- Tilt ile Dikey Vortisite Üretimi}

\[
\left(\frac{d\omega}{dt}\right)_{tilt} = \eta S
\]

\[
3\times 10^{-4} \, \text{s}^{-2}
\]

%%%%%%%%%%%%%%%%%%%%%%%%%%%%%%%%%%%%%%%%%%%%%%%%%%%%%%%%%%%%
\section*{Soru 11 -- Negatif Yüzdürme ve Rüzgâr}

\[
B = -0.25 \, \text{m/s}^2
\]

\[
\text{DCAPE} \approx 750\, \text{J/kg}
\]

\[
w_{\max} \approx 39\, \text{m/s}
\]

%%%%%%%%%%%%%%%%%%%%%%%%%%%%%%%%%%%%%%%%%%%%%%%%%%%%%%%%%%%%
\section*{Soru 12 -- Vortisten Basınç Farkı}

\[
\Delta p = \frac{1}{2} \rho \Omega^2 R^2
\]

\[
\Delta p \approx 8 \, \text{hPa}
\]

%%%%%%%%%%%%%%%%%%%%%%%%%%%%%%%%%%%%%%%%%%%%%%%%%%%%%%%%%%%%
\section*{Soru 13 -- Sağ Hareketli Süper Hücre}

Pozitif SRH → mesosiklon → düşük basınç → sağa sapma.

%%%%%%%%%%%%%%%%%%%%%%%%%%%%%%%%%%%%%%%%%%%%%%%%%%%%%%%%%%%%
\section*{Soru 14 -- Hodograf Eğriliği}

Saat yönünde eğri hodograf → sağ hareketli süper hücreyi güçlendirir.

%%%%%%%%%%%%%%%%%%%%%%%%%%%%%%%%%%%%%%%%%%%%%%%%%%%%%%%%%%%%
\section*{Soru 15 -- Kelvin Etkisi}

\[
e_s(r) = e_s(\infty) \exp\left(\frac{2\sigma}{\rho_w R_v T r}\right)
\]

Küçük damlada \(e_s\) büyüktür → buharlaşma eğilimi artar.

%%%%%%%%%%%%%%%%%%%%%%%%%%%%%%%%%%%%%%%%%%%%%%%%%%%%%%%%%%%%
\section*{Soru 16 -- Köhler Teorisi}

\[
\ln S = \frac{A}{r} - \frac{B}{r^3}
\]

\[
r_c = \left(\frac{3B}{A}\right)^{1/2}
\]

Minik NaCl taneciği → çok yüksek \(S_c\).

%%%%%%%%%%%%%%%%%%%%%%%%%%%%%%%%%%%%%%%%%%%%%%%%%%%%%%%%%%%%
\section*{Soru 17 -- Bergeron Süreci}

\[
e_{sw} > e_{si}
\]

Buz için süpersatürasyon ≈ %30.

%%%%%%%%%%%%%%%%%%%%%%%%%%%%%%%%%%%%%%%%%%%%%%%%%%%%%%%%%%%%
\section*{Soru 18 -- CAPE vs Dinamik Basınç}

CAPE updraft'ı başlatır;  
dinamik basınç (mesosiklon) olgun süper hücrede daha baskın olabilir.

\end{document}
